% !TeX root = courtade-kumar.tex
% !TEX TS-program = pdfLaTeX
\newpage
\label{sec:high-noise}
\section{The high noise inequality}
By negating $B$ as necessary, we can assume that
$\E[B]\le 1/2$. Let us observe that we can also assume
$\E[B]\ge1/6$, as otherwise the bound follows from
known inequalities.
\begin{lemma}
\label{lem:biasedB}
When $\E[B]\le1/6$ we have $\muti{B}{X}\le\muti{Y_1}{X}$.
\end{lemma}
\begin{proof}
Let us write
\begin{align*}
\muti{B}{X} &= \sum_i \muti{B}{X_i \emid X_{<i}}\\
  &\le \rho^2\sum_i \muti{B}{Y_i \emid X_{<i}}\\
  &\le \rho^2\sum_i \muti{B}{Y_i \emid Y_{<i}}\\
  &  = \rho^2 \muti{B}{Y}
\end{align*}
Setting $\mu=\E[B]$, we have $\muti{B}{Y}=\entb{\mu}$.
Direct calculation shows that when $\mu\le 1/6$ it holds
that $\entb{\mu}\le \frac{1}{2\ln 2}$. Observing
\begin{align*}
\frac{\rho^2}{2\ln 2} \le \kldivb{p}{1/2} = \muti{Y_1}{X}
\end{align*}
finishes the proof.
\end{proof}

\subsection{One-sided functions}
For $i\in[n]$ let us define $d_i =\E\sparen{B(2Y_i - 1)}$.
We show that if the output $B$ is highly dependent on a single
coordinate $Y_i$, an argument similar to 
\autoref{lem:biasedB} establishes the bound.

\begin{lemma}
\label{lem:high-corr}
Suppose $\E[B]>1/6$ and there exists a coordinate $i$
such that $\abs{d_i}> \frac{1}{6}$. 
Then we have $\muti{B}{X}\le\muti{C_\mu}{U}$.
\end{lemma}
\begin{proof}
Let us reveal $Y$ gradually as $X_i$, $Y_{-i}$ and
$Y_i$ and study the information they convey about $B$
using the chain rule:
\begin{align*}
\muti{B}{Y} &= \muti{B}{X_i} + \muti{B}{Y_{-i}\emid X_i}
  + \muti{B}{Y_i\emid Y_{-i}X_i}.
\end{align*}
Observe that $\muti{B}{Y}=\ent{B}$ as $Y$ determines $B$ and
$\muti{B}{Y_i\emid Y_{-i}X_i}\ge 2\abs{d_i} \ent{Y_i\emid X}$.
Using the chain rule again,
\begin{align*}
\muti{B}{X}
  &=   \muti{B}{X_i} + \muti{B}{X_{-i}\emid X_i}\\
  &\le \muti{B}{X_i} + \rho^2 \muti{B}{Y_{-i}\emid X_i}\\
  &\le \muti{B}{X_i} + \rho^2 \nparen{\ent{B}
    - \muti{B}{X_i}-d_i \ent{Y_i \emid X}}\\
  &=   (1-\rho^2)\muti{B}{X_i} + \rho^2 \ent{B}
    - 2\rho^2\abs{d_i} \ent{Y_i\emid X}
\end{align*}
Using \autoref{lem:scalarlemmathatidnthaveyet} we get
\begin{align*}
\muti{B}{X} \le \muti{C_\mu}{U},
\end{align*}
which completes the proof.
\end{proof}

\subsection{Even-handed functions}
\label{sec:well-rounded}
In the remainder of the paper, we show that even-handed
functions cannot come close to the benchmark $\muti{C_\mu}{U}$
due to Fourier theoretic reasons---any informative function
needs to be highly dependent on at least some coordinate.
In particular, assuming $1/6\le \E[B]\le1/2$ and
$\abs{d_i}\le 1/6$ for all $i\in[n]$ we show that
$\muti{B}{X}<\muti{C_\mu}{U}$.


Define
\begin{align*}
f(x) &= \E\sparen{B\emid Y=x}\text{, and}\\
g(x) &= \E\sparen{B\emid X=x}.
\end{align*}
Our proof proceeds as follows. We have 
\begin{align*}
\muti{B}{X} &= \kldiv{B\emid X}{B}\\
  &= \E_x\sparen{\kldivb{g(x)}{\mu}}
\end{align*}
We will design a function $Q_{\mu, p}(t)$ such that
$\kldivb{t}{\mu}\le Q_{\mu,p}(t)$ for all $t\in[0,1]$
further, $Q_{\mu,p}(t)$ is the minimum of two quadratic
polynomials in $t$.
The fact that $Q_{\mu,p}(t)$ is made of quadratic
polynomials in $t$ allows us to argue that
$\E_x[Q_{\mu,p}(g(x))]$ gets noticeably attenuated
via Fourier analysis. Moreover $Q_{\mu,p}$ is chosen
so that the attenuated value satisfies
\begin{equation*}
\E_x[Q_{\mu,p}(g(x))] \le \muti{C_\mu}{U},
\end{equation*}
which gives the desired bound.

Let us define the level energies
of $f$ as
\begin{align*}
W_k &= \sum_{\abs{S}=k}\E\sparen{\chi_S(Y)B}^2\\
    &= \sum_{\abs{S}=k}\E_x\sparen{\chi_S(x)f(x)}^2
\end{align*}
It is well known and not difficult to see that
$\E[\chi_S(X)B] = \rho^{|S|}\E[\chi_S(Y)B]$, since
$X$ is a noisy copy of $Y$. Our key observation is that
\begin{align*}
\E_x[g(x)^2] - \rho^2\E_x[f(x)g(x)]
  &=   \sum_{k}\nparen{\rho^{2k} - \rho^{k+2}} W_k\\
  &\le \rho^2(1-\rho) W_1 - \rho^2\mu^2.
\end{align*}

\begin{lemma}
Suppose $\abs{d_i}\le 1/6$ for all $i\in[n]$. Then $W_1\le 4/5$.
\end{lemma}
\begin{proof}
\end{proof}
